\subsection{Ninth Letter}

\begin{quotex}
Main points:

\begin{itemize}
\item Guenon gets a little snarky about the possible and the real 
\item Guenon considers Inayat Khan to be absurd, demonstrating that a ``legitimate" initiation offers no guarantee 
\item Guenon teaches Evola about Islam, including role of Melchizedek, the coming Imam, El-Khidr, and the meaning of infallibility and impeccability 
\item Guenon had high regard for \textbf{Mircea Eliade}. 
\item He tells a strange story about a spell cast on \textbf{Leon de Poncins}. We will have more to say about the latter at some other point. 
\item Guenon tells an even stranger story about a Jewish lawyer, but we are left hanging. 
\item Guenon explains why he never allows photos to be taken of him. 
\item Since Evola neither knew Guenon's age nor had ever seen a photograph of him, it proves that they had never met each other in person, no matter what others might claim. 
\end{itemize}
\end{quotex}

2 August 1949

Cairo, Egypt

[Discussions about publishing, translations, and the proofs of Revolt omitted]

Regarding \emph{Revolt}, you are doubtlessly correct: it will be simpler if you send me the proofs and I send back my observations after reading them, because otherwise you may have already modified something that I had noted.

\textbf{Inayat Khan}, whom I also knew, was regularly associated with the tariqah Chishtiya, one of the most widespread in India and totally orthodox; something that did not prevent the organization he founded to be completely the fruit of his fantasy and lacking any value; the name, ``order of Sufis" that he gave it is also truly absurd.

As for the Masonic question, I think like you that it is useless to revisit it again. I only point out that as you say yourself this time, it is good to understand how it is not possible to speak of ``Masonry" meaning a type of global entity, that in reality does not exist, or rather, if you prefer, exists only in the line of principle and which one cannot attribute it to any more or less exterior action: the refusal by some of its branches to recognize others whose deviations they criticize sufficiently proves on the other hand that there is no unity in this regard.

The problem of the possible and the real seems very simple and obvious to me, but, of course, under the condition of examining it from the metaphysical point of view. It is obvious that, from the philosophical point of view, one can always think anything whatsoever and discuss a problem endlessly without ever reaching a conclusion; it is even what characterizes profane speculation, and I have never been able to entertain any interest for those so-called ``problems" that fundamentally have only a verbal existence.

Melchizedek corresponds, in Islamic esoterism, to the function of the Qutb, as I have otherwise explained in \emph{King of the World}. On the other hand, El-Khider is the Master of the Afrad, which are found outside the jurisdiction of the Qutb and is said that they are not even known by it; in this regard, the Koranic story of the meeting between El-Khidr and Moses (Surat El-Kalif) is otherwise very significant. The way of the Afrad is something absolutely exceptional, and no one can choose it on his own initiative. It is about an initiation received beyond the ordinary means and belongs in reality to another chain (perhaps you can find an article of Abdul-Hadi in which he deals with these two chains, even if his definitions are not perhaps very clear).

In the Jewish Kabbalah, the same distinction is found expressed through the duality of Metatron and Sandalphon.

The hidden Imam is something completely different: those who admit his existence generally think that it is he who has to appear as the Mahdi. He is on the other hand defined as ``el-Muntazer", that could mean, the ``expected one", but that is interpreted almost always as ``he who waits".

Doctrinal infallibility belongs to whoever exercises legitimately a traditional function, naturally within the limits of that same function.

The issue of ``impeccability" is quite different, and it is usually considered, at least in orthodox tradition, as reserved to the Prophet: if it happens that he sometimes performs some acts that could seem reprehensible from an exterior point of view, it is only a question of an appearance, and such actions should in reality justify for the reason that elude the understanding of ordinary men.

I can provide yew some news about \textbf{Mircea Eliade}: he published, as you might know, three articles from his journal \emph{Zalmoxis}, the last of which in 1942. After that, he spent the rest of the war in Portugal, and subsequently he returned to Paris where he still is today. He has had many items published recently: \emph{Yoga: Immortality and Freedom}, the \emph{Myth of the Eternal Return}, and the \emph{History of Religious Ideas} (which I have not yet had time to read), without mentioning the many important articles in the \emph{Revue de l'Histoire des Religions}. I don't have his address, but I think I can easily find it out and will then let you know.

As for \textbf{Leon de Poncins}, it is a matter of a rather unpleasant story. Shortly before the war, a certain Eve Louguet was his secretary who took part in a group of dangerous sorcerers. He himself was a victim of these people and concerning the people who by chance saw him again around 1940, they reported to me that he seemed to have undergone a true collapse. I never knew what became of him since, but, in such conditions, I have many doubts that he can still be alive. [He actually outlived both Guenon and Evola. \flright{\textsc{ed.]}}

What is strange is that in the same period, one of the individuals in question tried to start a correspondence with me for some reason. At that time I did not know what it was about, but very soon the affair appeared suspicious to me, so that I immediately gave him a clean break.

I recently had the chance to speak about you with Mr. M., who has for more than a year been representing Argentina in Cairo and he informed me that he had known you at one time.

He intends to translate \emph{Man and his Becoming} into Spanish; up until now, only the \emph{General Introduction} has been translated into that language, in a version published in Buenos Aires during the war.

Since you asked me my age, I am 62 years old; I knew that you had to be younger than I, but I didn't think that the difference was so great. [Evola was 51 at the time.] As for my photograph, I am sorry that I cannot satisfy your request, but the truth is that I don't have any, and for many reasons. In fact, first of all what could be called a matter of principle that commits me, as you say, to give no importance to anything that is of a simply individual character. But, beyond that, I am also cautious that it could present some danger. About 15 years ago, I was informed that a certain Jewish lawyer was poking all around here to procure one of my photographs, claiming to be willing to pay any price. I never knew what he truly wanted to do with it, but what is certain in any case is that his intentions were not at all benevolent. Since one never knows where a photograph can end up, I concluded from the episode that it was much more prudent to not take one.



\flrightit{Posted on 2012-08-13 by Cologero }

\begin{center}* * *\end{center}

\begin{footnotesize}\begin{sffamily}



\texttt{Ali Alavi on 2012-08-13 at 14:32 said: }

Would I be wrong to suppose that the Abdul Hadi Guénon mentions and refers to, is none other than Ivan Aguéli?


\hfill

\texttt{Cologero on 2012-08-13 at 17:53 said: }

Ali Alavi, thank you for pointing that out as I neglected to include that in a note. Here is the note from the collection of letters:

\begin{quotex}
Abdul-Hadi, i.e., John Gustav Agelii (1869-1917), also called Ivan Aguéli, a painter, was the one who connected Rene Guenon to Islamic esoterism through the tariqa of sheik Ilaysh el-Kebir. 

\end{quotex}

\hfill

\texttt{Dominion on 2012-08-13 at 22:53 said: }

Not having read about Leon de Poncins before, does this account fit with what is known about his life before and after said period? And if this is his response to a single picture, I can only imagine how Guenon would have viewed Facebook and the internet…


\hfill

\texttt{Cologero on 2012-08-13 at 23:03 said: }

I did a google search, Dominion, on Leon de Poncins and Eva Louguet; Gornahoor showed up first, so it wasn't very helpful.

Leon de Poncins was a counter-revolutionary Catholic who saw Masonry and Judaism as hidden forces in Europe; actually not unlike Evola.

The notes point out that Leon lived another 25 years or so; hence, Guenon overestimated the negative effects. I have no idea why Leon's secretary would try psychic attacks against him.

As for the Internet, there may be safety in numbers. There are so many images, and not enough Jewish lawyers to use them for their own purposes.


\hfill

\texttt{Aghorable on 2012-08-14 at 10:07 said: }

Regarding how G. might view pervasive ``social networking", recall his describing a world where everything was made public as a ``monstrosity" (ROQ). Even then he had surmised enough to acknowledge this trend in the possibility of its full extension, and look around you now; yet another instance of a terminal approach to the limit in a particular domain (spearheaded by a nontraditional Jew).

For me, it necessitates a serious redefinition of where one invests oneself; the greater the encroachment, the more one is forced to redirect his essence, away from the glare of whatever fumbling hunchback's shoulders our world is currently resting on.

Concerning the psychic attacks from the secretary, it could have been a case of her mainly serving as a conduit for lower forces by her involvement in some peculiar group, without her necessarily knowing entirely what is going on. Damn, I bet she friend-zoned him too.

Generally, anyone with a practical interest in initiation does well to consider just what it would mean to a person living in, say, a typical metropolis of the modern West to suddenly gain a direct perception of their inner being, without the supports that existed in regularly organized dominions. Such an operation would imply, at some point, an unguarded opening to the psychic plane, and the lower forces that inevitably converge upon one are exceedingly difficult to detect as something extrinsic at first, perhaps because they are not entirely so. This is why G. mentions, more than once, the importance of theoretical preparation; should one face a collapse, he would need tools and strength to dig himself out and rebuild. And then he might find himself in a new state where he directly perceives the very forces that are behind the events, both petty and great, that are taking place every day in the corporeal realm, especially wherever he is situated. This is the `world as power'. and as you might imagine, living a life with such awareness is not for the faint hearted.

For the majority who have no idea of what is really at work when they are hit by forces of a similar kind, many end up in the hands of their local psychiatric team, which profession has apparently undergone a curious retreat from its previous position at the avant-garde of the psychic incursion, to the reassuring illusions of security in materialism and its favourite panacea, the chemical dose. Still others end up dead, often in grotesque circumstances (search your local Coroner's reports for particulars), hence G.'s doubts about the Vicomte de Poncins's survival.


\hfill

\texttt{Cologero on 2012-08-14 at 22:08 said: }

I haven't found anything about Poncin's personal life other than he was from an old aristocratic family and was a Catholic counter-revolutionary along the lines of Joseph de Maistre and Donoso Cortes. I have some old books that reference his works and confirm that point of view. Furthermore, in the Italian edition of Evola's Revolt, Poncin's work, the Hidden (or Occult) War, is referred to twice. The references are not included in the English translation, nor are the references to Donoso. Evola himself described Poncin as a Traditionalist.


\hfill

\texttt{Lucio on 2016-06-28 at 21:28 said: }

``I recently had the chance to speak about you with Mr. M., who has for a time more than a year been representing Argentina in Cairo and he informed me that he had known you at one time".

Who is this Mr. M.?


\hfill

\texttt{Cologero on 2016-06-28 at 22:55 said: }

Lucio: I assume it was a Maxim Montoto who translated ``Man and his Becoming" into Spanish (as Guenon indicated). However, I can't find out anything about him.


\hfill

\texttt{Lucio on 2016-06-29 at 02:33 said: }

Thanks, I consider important to know who the translators are, and I assume that the edition I have of ``General Introduction" is from him, though I cannot tell for sure.


\end{sffamily}\end{footnotesize}
